\section{Introduction}

Nearest-neighbor search returns the vectors with the smallest distances or the
largest similarities to a query. In flat search, the system computes $N$ scores
for each query and then selects the best $K$ entries. FAISS showed that this
$K$-selection step can be a substantial part of the runtime after the distance
kernel is optimized \citep{johnson2019billion,douze2024faiss}.

This paper studies a narrow but common setting: after a flat-search score row is
computed, the remaining task is to select its top-$K$ entries. Monte-Carlo TopK
is a simple selector for this setting. It samples a small set of scores, uses a
sampled order statistic as a threshold, scans the full score row once, and runs
exact top-$K$ on the candidates below the threshold. The method does not reduce
the full distance computation, change the index, or change the vector representation;
Section~\ref{sec:method} defines this selector and its exact-mode checks. It is
useful when exact selection is expensive relative to the rest of the pipeline.

The main practical issue is the threshold. If it keeps fewer than $K$
candidates, the row cannot return $K$ results. If it keeps too many candidates,
the candidate buffer overflows. The exact mode handles both cases by falling
back to full-row exact selection. A separate capped mode avoids overflow
fallbacks by keeping only a bounded number of candidates, but this mode is
approximate. The experiments therefore report speed, candidate-count tails,
fallback counts, returned-score agreement, and identifier overlap separately.
This separation matters because equal scores at the $K$-th boundary can make two
valid exact selectors return different identifiers.

The FAISS GPU results show that the approach can give a large selector gain in
flat squared-L2 search. On four ANN10M datasets, Monte-Carlo TopK gives a 15.0x
selector speedup and a 2.62--2.83x end-to-end speedup with zero fallbacks across
five sample seeds. On SIFT1M, a no-fallback setting gives a 10.0x selector
speedup and a 2.72x end-to-end speedup. These gains are not from reducing
distance work; they come from reducing the cost of the final selection.

The contributions are:

\begin{enumerate}
  \item A post-score top-$K$ selector that uses a sampled threshold, a candidate
        buffer, and exact top-$K$ selection over candidates.
  \item An exact-mode contract with explicit underflow and overflow fallbacks,
        plus a separate approximate capped mode.
  \item FAISS CPU and GPU experiments on SIFT1M and four ANN10M datasets,
        including speed breakdowns, candidate tails, fallback counts, and
        tie-aware correctness checks.
\end{enumerate}
